% ID: MAT-TRI-EQG-011
% Titolo: Dominio e zeri di una funzione goniometrica (2)
% Disciplina: Matematica
% Area: Trigonometria
% Argomento: Equazioni goniometriche
% Classe: Quarta secondaria di secondo grado
% Difficolta: 4
% Tipo: problema_numerico
% Tag: matematica, trigonometria, equazioni_goniometriche
% Fonte: verifica_goniometria_intermedia_anonimizzata
% Provenienza: materiale fornito dal docente; testo anonimizzato; soluzione redatta per Exergo

\begin{esercizio}
Determina il dominio e gli zeri della funzione: 
\[
f(x)=\dfrac{\sin 2x-\cos x}{\sin x\,(2\cos x-\sqrt{3})}.
\]
Salvo diversa indicazione, lavora in $\mathbb R$; angoli in radianti.
\end{esercizio}

\begin{soluzione}
Il denominatore impone $x\ne k\pi$ e $x\ne\pm\pi/6+2k\pi$.
Il numeratore è $\cos x(2\sin x-1)$.
Gli zeri ammessi sono $x=\pi/2+k\pi$ e $x=5\pi/6+2k\pi$;
$\pi/6+2k\pi$ è escluso. In tutte le famiglie di soluzioni $k\in\mathbb{Z}$.
\end{soluzione}
